% !TEX root = ../main.tex

\section{Conclusion}
Recall the reaction equation:
$$\rxn$$

In test tubes B and C, when we added \ce{KSCN} and \ce{FeCl3} respectively, we observed a shift towards the products. This is supported by the observed colour change to a redder solution. When the concentration of the reactants was increased, the forward reaction rate increased due to more frequent collisions. The reverse reaction rate remained unchanged as the reverse reaction does not require the individual species to carry forward. More \textit{red} products were being formed from \textit{yellow} reactants, and the system shifted to the right to reach a new equilibrium.

% \\~\\

For test tube E, consider the reaction:
$$\ce{Fe^3+ + 3OH- -> Fe(OH)3_{(s)}}$$
where Iron reacts with Hydroxide to form a precipitate.

% \\~\\

When we added \ce{NaOH}, we observed a shift towards the reactants. This was supported by the observed colour change to a yellower solution. When the \ce{NaOH} was added, the \ce{OH-} ions reacted with the \ce{Fe^3+} ions to form a precipitate. This resulted in a decrease in the concentration of \ce{Fe^3+}. The decrease in concentration resulted in a decrease in collisions, which in turn decreased the forward rate. As the reverse rate does not depend on the concentration of the reactants, it remained unchanged. More reactants were being produced than products, and the system shifted to the left to reach a new equilibrium.

For test tubes A and D, we observed no shift in the equilibrium. This is supported by the colour of the solution remaining unchanged. For test tube A, no chemicals were added, meaning there was no stress in the system. It follows that there was then no shift in the equilibrium. For test tube D, adding \ce3{KCl} did not affect the equilibrium as the \ce{K+} and \ce{Cl-} ions do not participate in the reaction. This results in no stress in the system, and thus no shift in the equilibrium.